% GENERATED FILE. Edit canonical lessons, then run npm run generate:books.
% canonical-source-hash: fnv1a64:4e2b2c52b858041d
% canonical-lessons: KA-C305-arcaka, KA-C305-kalavida, KA-C305-gayaka, KA-C305-nata, KA-C305-patrakarta

\chapter{Temple priest, Artist, Singer, Actor, Journalist}
\label{ch:ka-temple-priest-artist-singer-actor-journalist}
\hlchaptermodality{305}

\begin{chapteropening}
\textbf{By the end of this chapter:} \emph{I can say a temple priest, an artist, a singer, an actor and a journalist.}

Complete the last lesson of chapter 305: I can say a temple priest, an artist, a singer, an actor and a journalist.
\end{chapteropening}

\section[\texorpdfstring{\kn{ಅರ್ಚಕ} (arcaka)}{ಅರ್ಚಕ (arcaka)}]{\kn{ಅರ್ಚಕ} (arcaka) — a temple priest}

\label{lesson:KA-C305-arcaka}



\begin{quote}
Before the new one: say the Kannada for a subject, then the Kannada for mathematics.
\end{quote}

\subsection*{You'll want to know: \kn{ಅರ್ಚಕ}}
\textbf{\kn{ಅರ್ಚಕ}} — \emph{arcaka} — \textquotedblleft{}a temple priest\textquotedblright{}.

\begin{grammarlens}[title={What you've built}]
One more word for this chapter.
\end{grammarlens}

\subsection*{Your turn}
\begin{itemize}
\raggedright
  \item \emph{Say it:} \emph{arcaka}
  \item \emph{Say it:} \emph{arcaka}, once more
  \item \emph{Say it:} say \emph{arcaka}
  \item \emph{Recall:} say the Kannada for a subject, then the Kannada for mathematics, then say \emph{arcaka} again
\end{itemize}

\begin{tcolorbox}[breakable,colback=teal!4,colframe=teal!35!black,title={Before you move on}]
What does \kn{ಅರ್ಚಕ} mean? (\textquotedblleft{}A temple priest\textquotedblright{}.) Say it once more.
\end{tcolorbox}
\section[\texorpdfstring{\kn{ಕಲಾವಿದ} (kalāvida)}{ಕಲಾವಿದ (kalāvida)}]{\kn{ಕಲಾವಿದ} (kalāvida) — an artist}

\label{lesson:KA-C305-kalavida}



\begin{quote}
Before the new one: say the Kannada for mathematics, then the Kannada for a temple priest.
\end{quote}

\subsection*{You'll want to know: \kn{ಕಲಾವಿದ}}
\textbf{\kn{ಕಲಾವಿದ}} — \emph{kalāvida} — \textquotedblleft{}an artist\textquotedblright{}.

\begin{grammarlens}[title={What you've built}]
One more word for this chapter.
\end{grammarlens}

\subsection*{Your turn}
\begin{itemize}
\raggedright
  \item \emph{Say it:} \emph{kalāvida}
  \item \emph{Say it:} \emph{kalāvida}, once more
  \item \emph{Say it:} say \emph{kalāvida}
  \item \emph{Recall:} say the Kannada for mathematics, then the Kannada for a temple priest, then say \emph{kalāvida} again
\end{itemize}

\begin{tcolorbox}[breakable,colback=teal!4,colframe=teal!35!black,title={Before you move on}]
What does \kn{ಕಲಾವಿದ} mean? (\textquotedblleft{}An artist\textquotedblright{}.) Say it once more.
\end{tcolorbox}
\section[\texorpdfstring{\kn{ಗಾಯಕ} (gāyaka)}{ಗಾಯಕ (gāyaka)}]{\kn{ಗಾಯಕ} (gāyaka) — a singer}

\label{lesson:KA-C305-gayaka}



\begin{quote}
Before the new one: say the Kannada for a temple priest, then the Kannada for an artist.
\end{quote}

\subsection*{You'll want to know: \kn{ಗಾಯಕ}}
\textbf{\kn{ಗಾಯಕ}} — \emph{gāyaka} — \textquotedblleft{}a singer\textquotedblright{}.

\begin{grammarlens}[title={What you've built}]
One more word for this chapter.
\end{grammarlens}

\subsection*{Your turn}
\begin{itemize}
\raggedright
  \item \emph{Say it:} \emph{gāyaka}
  \item \emph{Say it:} \emph{gāyaka}, once more
  \item \emph{Say it:} say \emph{gāyaka}
  \item \emph{Recall:} say the Kannada for a temple priest, then the Kannada for an artist, then say \emph{gāyaka} again
\end{itemize}

\begin{tcolorbox}[breakable,colback=teal!4,colframe=teal!35!black,title={Before you move on}]
What does \kn{ಗಾಯಕ} mean? (\textquotedblleft{}A singer\textquotedblright{}.) Say it once more.
\end{tcolorbox}
\section[\texorpdfstring{\kn{ನಟ} (naṭa)}{ನಟ (naṭa)}]{\kn{ನಟ} (naṭa) — an actor}

\label{lesson:KA-C305-nata}



\begin{quote}
Before the new one: say the Kannada for an artist, then the Kannada for a singer.
\end{quote}

\subsection*{You'll want to know: \kn{ನಟ}}
\textbf{\kn{ನಟ}} — \emph{naṭa} — \textquotedblleft{}an actor\textquotedblright{}.

\begin{grammarlens}[title={What you've built}]
One more word for this chapter.
\end{grammarlens}

\subsection*{Your turn}
\begin{itemize}
\raggedright
  \item \emph{Say it:} \emph{naṭa}
  \item \emph{Say it:} \emph{naṭa}, once more
  \item \emph{Say it:} say \emph{naṭa}
  \item \emph{Recall:} say the Kannada for an artist, then the Kannada for a singer, then say \emph{naṭa} again
\end{itemize}

\begin{tcolorbox}[breakable,colback=teal!4,colframe=teal!35!black,title={Before you move on}]
What does \kn{ನಟ} mean? (\textquotedblleft{}An actor\textquotedblright{}.) Say it once more.
\end{tcolorbox}
\section[\texorpdfstring{\kn{ಪತ್ರಕರ್ತ} (patrakarta)}{ಪತ್ರಕರ್ತ (patrakarta)}]{\kn{ಪತ್ರಕರ್ತ} (patrakarta) — a journalist}

\label{lesson:KA-C305-patrakarta}



\begin{quote}
Before the new one: say the Kannada for a singer, then the Kannada for an actor.
\end{quote}

\subsection*{You'll want to know: \kn{ಪತ್ರಕರ್ತ}}
\textbf{\kn{ಪತ್ರಕರ್ತ}} — \emph{patrakarta} — \textquotedblleft{}a journalist\textquotedblright{}.

\begin{grammarlens}[title={What you've built}]
One more word for this chapter.
\end{grammarlens}

\subsection*{Your turn}
\begin{itemize}
\raggedright
  \item \emph{Say it:} \emph{patrakarta}
  \item \emph{Say it:} \emph{patrakarta}, once more
  \item \emph{Say it:} say \emph{patrakarta}
  \item \emph{Recall:} say the Kannada for a singer, then the Kannada for an actor, then say \emph{patrakarta} again
\end{itemize}

\begin{tcolorbox}[breakable,colback=teal!4,colframe=teal!35!black,title={Before you move on}]
What does \kn{ಪತ್ರಕರ್ತ} mean? (\textquotedblleft{}A journalist\textquotedblright{}.) Say it once more.
\end{tcolorbox}
